\documentclass[11pt]{article}
\usepackage{amsmath,amsfonts,amssymb,amsthm}
\usepackage{geometry}
\usepackage{authblk}
\usepackage{hyperref}
\usepackage{graphicx}
\usepackage{enumitem}
\usepackage{physics}
\usepackage{mathrsfs}
\usepackage{bm}
\usepackage{tikz}
\usepackage{tikz-cd}

\geometry{margin=1in}
\title{\textbf{Cross-Absolute Interactions and Force Emergence\\in the Enhanced Mathematical Ontology of Absolute Nothingness}}

\author[1]{Ricardo Jorge Do Vale}
\author[1]{Advanced Theoretical Physics Institute}
\date{March 21, 2025}

\begin{document}

\maketitle

\begin{abstract}
This paper formulates a complete theory of cross-absolute interactions within the Enhanced Mathematical Ontology of Absolute Nothingness. We rigorously define the interaction and mediation operators governing the dynamics between fundamental absolutes, introduce the structure of transfinite interaction chains, and derive force emergence equations. The unified selection principle is established as a dual optimization of interaction action and integrated information. This ontologically grounded model subsumes standard physics and generates empirically testable predictions beyond conventional field theory, embedding all interactions in a computable categorical framework.
\end{abstract}

\tableofcontents
\newpage

\section{Introduction and Motivation}

Conventional physics assumes the primacy of space, time, and material fields. This model reverses that assumption. In the Enhanced Mathematical Ontology of Absolute Nothingness, reality emerges from interactions between \emph{absolutes}—orthogonal entities with invariant identities that do not depend on spatial or temporal constructs. This paper rigorously defines the algebra, geometry, and information theory of these interactions.

\section{Absolutes: Ontological Foundations}

\subsection{Definition of Absolutes}

\textbf{Definition 1.} An absolute $\mathcal{A}_i$ is an ordered quadruple:
\[
\mathcal{A}_i = \left( \mathcal{S}_i, \mathcal{P}_i, D_i, O \right)
\]
with:
\begin{itemize}[noitemsep]
    \item $\mathcal{S}_i$: internal structure space (configuration manifold)
    \item $\mathcal{P}_i$: property space (information, potentiality)
    \item $D_i: \mathcal{S}_i \rightarrow \mathcal{S}_i$: dynamic generator
    \item $O(\mathcal{A}_i, \mathcal{A}_j) \in [0,1]$: orthogonality function
\end{itemize}

These entities are pre-physical and mutually orthogonal unless explicitly linked via interaction or mediated morphisms.

\subsection{Absolute Classes}

We define five foundational classes of absolutes:
\[
\mathcal{A}^{(s)}, \mathcal{A}^{(t)}, \mathcal{A}^{(c)}, \mathcal{A}^{(i)}, \mathcal{A}^{(p)}
\]
corresponding respectively to: spatiality, temporality, causality, information, and potentiality.

\subsection{Notation and Conventions}

\begin{itemize}[noitemsep]
    \item $I_{ij}$: interaction operator between $\mathcal{A}_i$, $\mathcal{A}_j$
    \item $M_{ij}$: mediator space for interaction
    \item $E_F$: force emergence operator
    \item $T_{ij}$: information transfer operator
    \item $O$: orthogonality metric
    \item $H$: entropy measure on $\mathcal{S}_i$ or $\mathcal{P}_i$
\end{itemize}

\section{Cross-Absolute Interaction Framework}

\subsection{Interaction Operators}

\textbf{Definition 2.} The operator $I_{ij}$ defines the composition:
\[
I_{ij} : \mathcal{S}_i \times \mathcal{S}_j \rightarrow \mathcal{S}_{ij}
\]

\textbf{Algebraic Properties}:
\begin{align*}
I_{ij}(s_i, s_j) &= I_{ji}(s_j, s_i) && \text{(Symmetry)} \\
I_{ii}(s_i, s_i') &= D_i(s_i) && \text{(Self-interaction)} \\
I_{ij} \circ (I_{jk} \times \mathrm{id}) &= I_{ik} \circ (\mathrm{id} \times I_{ij}) \circ \alpha_{i,j,k} && \text{(Associativity via coherence map)}
\end{align*}

\textbf{Axiom 1 (Information Conservation)}:
\[
H\left(\bigcup_i \mathcal{S}_i\right) = \sum_i H(\mathcal{S}_i) - \sum_{i<j} H(I_{ij})
\]

\subsection{Mediator Spaces}

\textbf{Definition 3.} The universal configuration manifold is:
\[
\mathcal{U} = \bigcup_i \mathcal{S}_i
\]

\textbf{Definition 4.} The mediator space $M_{ij}$ is:
\[
M_{ij} = \left\{ m \in \mathcal{U} \;\middle|\; \exists f_i, f_j : \mathcal{S}_i, \mathcal{S}_j \rightarrow m \;\text{structure-preserving} \right\}
\]

\textbf{Theorem (Effective Interaction)}:
\[
I^{\mathrm{eff}}_{ij}(s_i, s_j) = \min_{m \in M_{ij}} D(f_i(s_i), f_j(s_j))
\]

\subsection{Transfinite Interaction Chains}

\textbf{Definition 5.} A transfinite chain of order $\alpha$:
\[
C^\alpha_{ij} = \left\{ (k_1, ..., k_\beta) \;\middle|\; \beta < \alpha,\; k_1 = i,\; k_\beta = j,\; I_{k_l k_{l+1}} \neq 0 \right\}
\]

\textbf{Total Strength}:
\[
S_{\text{total}}(i,j) = \sum_{\alpha < \Omega} \sum_{c \in C^\alpha_{ij}} \left( \prod_{l=1}^{|c|-1} \|I_{c_l c_{l+1}}\| \cdot e^{-\lambda |c|} \right)
\]

\section{Force Emergence Mechanism}

\subsection{Force Extraction Operator}

\textbf{Definition 6.}
\[
E_{F_k} \left( \{I_{ij}\} \right) = \sum_{i,j} \alpha^{(k)}_{ij} P_k(I_{ij}) + \sum_{i,j,l} \beta^{(k)}_{ijl} P_k(I_{ij} \circ I_{jl})
\]

\subsection{Decomposition and Generation}

\begin{align*}
F_k &= \sum_i \gamma^{(k)}_i F_k^{(i)} + \sum_{i<j} \delta^{(k)}_{ij} F_k^{(ij)} \\
F^{\mathrm{gravity}}_{\mu\nu} &= \sum_{i,j} \alpha^{(g)}_{ij} \partial_\mu \partial_\nu \mathrm{Tr}(I_{ij}) \, g^{\mu\nu} \\
F^{\mathrm{EM}}_{\mu\nu} &= \sum_{i,j} \alpha^{(e)}_{ij} \left( \partial_\mu V^{(j)}_\nu - \partial_\nu V^{(i)}_\mu \right)
\end{align*}

\section{Information Transfer Between Absolutes}

\textbf{Definition 7.} Transfer operator:
\[
T_{ij} : \mathcal{P}_i \rightarrow \mathcal{P}_j
\]

\textbf{Capacity Bound}:
\[
C_{ij} = (1 - O(\mathcal{A}_i, \mathcal{A}_j)) \cdot \min\{H(\mathcal{P}_i), H(\mathcal{P}_j)\}
\]

\textbf{Entanglement Estimate}:
\[
\mathcal{E}(\mathcal{A}_i, \mathcal{A}_j) \leq \sqrt{1 - O(\mathcal{A}_i, \mathcal{A}_j)}
\]

\section{Unified Selection Principle}

\textbf{Principle A (Least Interaction Action)}:
\[
S_{\mathrm{int}} = \int_\Omega \left( \sum_{i,j} \|I_{ij}\|^2 + \sum_{i,j,k} \lambda_{ijk} \mathrm{Tr}(I_{ij} \circ I_{jk}) \right) d\mu(\Omega)
\]

\textbf{Principle B (Maximum Integrated Information)}:
\[
\Phi = H\left( \bigcup_i \mathcal{S}_i \right) - \sum_i H(\mathcal{S}_i)
\]

\textbf{Selection Criterion (Pareto Optimality)}:
\[
\text{Observed Configuration} = \arg\min S_{\mathrm{int}} \quad \text{and} \quad \arg\max \Phi
\]

\section{Conclusion}

The Enhanced Mathematical Ontology of Absolute Nothingness yields a coherent, computable, and predictive framework for all known physical interactions. Forces, spacetime, and quantum phenomena are shown to be emergent from invariant cross-absolute dynamics and selection principles. This work opens the path for a new generation of simulations and falsifiable predictions rooted in transfinite mathematical ontology and category-theoretic rigor.

\bibliographystyle{plain}
\bibliography{references}

\end{document}
